\chapter{Knowledge and Skills Evaluation}
\label{chap:knowledge-skills-evaluation}

This chapter evaluates the knowledge and skills demonstrated through
the Vid-Scratch project by comparing the project evidence with the
expected curriculum rubric levels. The evaluation considers five main
dimensions: Software Development and Programming, Machine Learning /
AI Systems, Data Management / Knowledge Systems, Critical and Creative
Thinking, and Collaboration and Communication.

The evaluation is based on evidence from the project report, including
the implemented desktop application, AI pipeline, requirements
analysis, architecture diagrams, user interface workflow, experiments,
limitations, and project documentation. Although the original rubric
also includes Lifelong Learning, this chapter focuses on the five
dimensions used in the radar chart and summary table.

\section{Comparison Between Curriculum Targets and AI-Based Evaluation}
\label{section:curriculum-ai-comparison}

\subsection{Radar Chart Comparing Expected and Assessed Levels}
\label{subsection:rubric-radar-chart}

Figure~\ref{fig:rubric-radar-chart} compares the curriculum expected
level with the AI-assessed level for each evaluated dimension. The
chart shows that Vid-Scratch meets the expected level in Software
Development, Machine Learning / AI Systems, and Critical and Creative
Thinking. The main gaps appear in Data Management / Knowledge Systems
and Collaboration and Communication.

\begin{figure}[H]
    \centering
    \includegraphics[width=0.92\textwidth]{chapter_9/rubric_evaluation_radar.png}
    \caption{Rubric evaluation radar chart comparing curriculum expected levels and AI-assessed levels.}
    \label{fig:rubric-radar-chart}
\end{figure}

\subsection{Summary of Evaluation Results}
\label{subsection:evaluation-summary-table}

The detailed evaluation results are summarized in
Table~\ref{tab:rubric-evaluation-summary}. The assessment uses a
four-level rubric, where higher levels indicate stronger evidence of
technical maturity, systematic practice, and project-level impact.

\begin{table}[H]
\centering
\caption{Summary of curriculum expected levels and AI-assessed levels.}
\label{tab:rubric-evaluation-summary}
\renewcommand{\arraystretch}{1.3}
\begin{tabularx}{\textwidth}{|X|c|c|X|}
\hline
\textbf{Dimension} & \textbf{Expected} & \textbf{Assessed} & \textbf{Final Judgment} \\
\hline
Software Development and Programming & 2 & 2 & Meets expectation \\
\hline
Machine Learning / AI Systems & 3 & 3 & Meets expectation \\
\hline
Data Management / Knowledge Systems & 2 & 1 & Below expectation \\
\hline
Critical and Creative Thinking & 3 & 3 & Meets expectation \\
\hline
Collaboration and Communication & 3 & 2 & Below expectation \\
\hline
\end{tabularx}
\end{table}

Overall, the Vid-Scratch report demonstrates a strong undergraduate
software engineering and AI project. The strongest areas are Machine
Learning / AI Systems and Critical and Creative Thinking, where the
project provides clear evidence of adversarial video generation,
action-recognition robustness testing, evaluation-driven
experimentation, and honest discussion of limitations. The weaker
areas are Data Management / Knowledge Systems and Collaboration and
Communication, where the report provides less evidence of dataset
governance, versioning, access control, issue tracking, pull request
workflow, and team-level decision logs.

\section{Detailed Evaluation}
\label{section:detailed-evaluation}

\subsection{Software Development and Programming}
\label{subsection:software-development-evaluation}

The AI-assessed level for Software Development and Programming is
Level 2, which meets the expected curriculum level. The project
provides evidence of a complete application with a clear software
stack, user interface workflow, backend pipeline, version control, and
testing.

Vid-Scratch is implemented as a desktop application using Tauri, with
a React and TypeScript frontend and a Python backend pipeline. The
architecture includes a frontend layer, Tauri command layer, and Python
processing engine, showing structured integration across multiple
software components. The report also includes domain, class, and
sequence diagrams that describe the major system components and their
interactions.

The user workflow covers video selection, configuration, processing,
progress monitoring, output preview, and export. In addition, the test
section includes an inference-time benchmark, a training-time poisoning
experiment, and unit tests for core pipeline modules. Git and GitHub
are also used for version control.

The project does not reach Level 3 because the report does not provide
strong evidence of scalable deployment, containerization, cloud
orchestration, mature release engineering, regression testing strategy,
or formal DevOps practice. However, it clearly exceeds Level 1 because
it is a usable desktop application rather than a simple script.

\subsection{Machine Learning / AI Systems}
\label{subsection:ml-ai-evaluation}

The AI-assessed level for Machine Learning / AI Systems is Level 3,
which meets the expected curriculum level. Although Vid-Scratch does
not implement an LLM or agentic system, it demonstrates advanced
AI-system design through model-targeted adversarial optimization and
quantitative evaluation.

The project targets SlowFast R101, a pretrained temporal
action-recognition model, and evaluates adversarial vulnerability under
both inference-time and training-time scenarios. The pipeline combines
Bayesian Optimisation for keyframe selection, learned spatial
flow-field transformation, and gradient-based additive adversarial
noise. It also uses SSIM-based perceptual similarity constraints to
balance attack effectiveness with visual quality.

The report defines measurable AI task requirements, including
misclassification at inference time, degradation during training,
visual similarity preservation, and identification of important
frames. Evaluation metrics include attack success rate, classification
behavior, loss trends, and perceptual similarity. The benchmark also
reloads the saved poisoned video and reruns inference to confirm that
the attack remains effective after codec compression.

The project does not reach Level 4 because it does not present a new
publishable AI method or enterprise-level AI strategy. However, it is
more advanced than Level 2 because it includes model-targeted
optimization, threat-scenario framing, and metric-based evaluation.

\subsection{Data Management / Knowledge Systems}
\label{subsection:data-management-evaluation}

The AI-assessed level for Data Management / Knowledge Systems is
Level 1, which is below the expected curriculum level of Level 2. The
project uses datasets and video files, but it does not provide enough
evidence of systematic dataset management, metadata governance,
versioning, lineage tracking, access control, or knowledge-system
design.

The report states that data was sourced from platforms such as Kaggle,
Pexels, and Internet Archive, and the system processes MP4 videos by
loading frames with OpenCV. This shows basic unstructured-data
handling and dataset use. However, the report does not describe a
formal dataset preparation pipeline, metadata schema, dataset
versioning, cloud storage policy, or access-control mechanism.

For this reason, the project fits Level 1 rather than Level 2. To
reach Level 2, the report should include clearer evidence of dataset
organization, versioned data preparation, reproducible file structure,
metadata tracking, and possibly experiment data logs.

\subsection{Critical and Creative Thinking}
\label{subsection:critical-thinking-evaluation}

The AI-assessed level for Critical and Creative Thinking is Level 3,
which meets the expected curriculum level. The report demonstrates
systematic reasoning through problem framing, literature comparison,
threat-scenario definition, measurable evaluation, iteration,
limitations, and future work.

The project separates the problem into inference-time and training-time
attack scenarios. The literature review compares related approaches
such as DeepSAVA, BTC-UAP, stAdv, ART, and Foolbox, then identifies a
gap between research-oriented adversarial tools and user-facing
software. The project also defines measurable goals using attack
success rate, accuracy or loss behavior, and SSIM.

The report discusses trade-offs between attack effectiveness,
perceptual similarity, GPU cost, target model dependency, and short
video constraints. It also documents limitations and future work,
including GPU requirement, limited training-time impact, fixed target
model, short-video constraint, transformer-based models, and
GPU-efficient frame selection.

The project does not reach Level 4 because it does not create a new
organizational standard or broader methodology beyond the project
context. However, it exceeds Level 2 because it uses systematic
testing, metrics, risk analysis, and iterative evaluation.

\subsection{Collaboration and Communication}
\label{subsection:collaboration-evaluation}

The AI-assessed level for Collaboration and Communication is Level 2,
which is below the expected curriculum level of Level 3. The report is
well documented and communicates system constraints, risks,
stakeholders, design, and evaluation results clearly, but it does not
provide strong evidence of Level 3 collaboration practice.

The report includes background, requirements, architecture, AI
components, development process, deliverables, testing, limitations,
and future work. It identifies stakeholder groups such as general
users, students, researchers, educators, developers, ML practitioners,
and supervisors. It also communicates important project constraints,
including GPU dependency, pretrained model dependency, limited
training-time attack impact, fixed target model, and short-video
constraint.

Git and GitHub are used for version control, and the project was
reviewed through advisor feedback and project defense. However, the
report does not provide enough concrete evidence of GitHub Issues,
pull requests, traceable decision logs, cross-disciplinary leadership,
or formal team coordination. Therefore, Level 2 is a fair assessment.

\section{Gap Analysis}
\label{section:gap-analysis}

Table~\ref{tab:gap-analysis} summarizes the gap between the curriculum
expected level and the AI-assessed level.

\begin{table}[H]
\centering
\caption{Gap analysis between expected and assessed rubric levels.}
\label{tab:gap-analysis}
\renewcommand{\arraystretch}{1.3}
\begin{tabularx}{\textwidth}{|X|c|c|c|X|}
\hline
\textbf{Dimension} & \textbf{Expected} & \textbf{Assessed} & \textbf{Gap} & \textbf{Status} \\
\hline
Software Development and Programming & 2 & 2 & 0 & Meets expectation \\
\hline
Machine Learning / AI Systems & 3 & 3 & 0 & Meets expectation \\
\hline
Data Management / Knowledge Systems & 2 & 1 & -1 & Below expectation \\
\hline
Critical and Creative Thinking & 3 & 3 & 0 & Meets expectation \\
\hline
Collaboration and Communication & 3 & 2 & -1 & Below expectation \\
\hline
\end{tabularx}
\end{table}

The strongest dimensions are Machine Learning / AI Systems and
Critical and Creative Thinking. These areas are supported by the
project's adversarial machine learning pipeline, quantitative
evaluation, research-based method selection, and clear discussion of
technical limitations. The main gaps are in Data Management and
Collaboration and Communication. These gaps are mainly caused by
limited evidence of dataset governance, versioning, access control,
GitHub issue or pull request workflow, and team-level decision logs.

\section{Final Evaluation Conclusion}
\label{section:evaluation-conclusion}

The Vid-Scratch project shows that the student is capable of building
a meaningful AI-enabled software system that combines desktop
application development, adversarial machine learning, video
processing, model evaluation, and academic reporting. The project is
especially strong in its technical AI pipeline, clear problem framing,
measurable evaluation strategy, and honest discussion of limitations.

The student appears ready for undergraduate-level software engineering
and applied AI project work, especially in areas related to AI
application development, computer vision experimentation, and
robustness testing. To better meet the full program outcomes, the
project should strengthen evidence in data-management practice and
collaboration process. In future work, this could be improved by
documenting dataset versioning, experiment tracking, access control,
GitHub Issues and pull requests, continuous integration evidence, and
decision logs.
